\chapter{Option Contracts and the Options Exchanges}\label{ch:m1:option-contracts-and-exchanges}

A large American company has one share. On that share, this morning, there
are some two thousand different option contracts: calls and puts, at seventy
\omterm{def:m1:option-contracts-and-exchanges:option}{strike prices}, for expiries from this Friday to two years out. Each of the
two thousand is quoted, two-sided, on more than fifteen exchanges at once,
by \omterm{def:m1:what-a-trading-firm-does:market-maker}{market makers} who must update every one of those quotes whenever the share
moves a cent. The equity market of \cref{ch:m1:us-equity-market-structure}
has a few thousand instruments; the options market built on top of it has
hundreds of times as many. This chapter describes the contract, the single
clearing house that makes all those exchanges one market, and the evening of
the third Friday, when options turn back into shares. What an option is
\emph{worth} is the subject of One Quant Books 5 and 6; here the question is
what it \emph{is}.

\section{The contract}

\begin{definition}[Call, put, strike, expiry, premium]\label{def:m1:option-contracts-and-exchanges:option}
A \emph{call option}\index{call option} gives its holder the right, and not
the obligation, to buy the underlying at a fixed \emph{strike
price}\index{strike price} on or before an \emph{expiry}\index{expiry} date;
a \emph{put option}\index{put option} gives the right to sell. The price
paid for the option is its \emph{option premium}\index{option premium}. The
seller, or writer, receives the premium and bears the obligation.
\end{definition}

\begin{definition}[Multiplier, series, moneyness]\label{def:m1:option-contracts-and-exchanges:series}
The \emph{option multiplier}\index{option multiplier} is the quantity of
underlying per contract: 100 shares for a standard US equity option, \$100
per \omterm{def:m1:index-and-single-stock-futures-worldwide:point}{index point} for the large index options. An \emph{option
series}\index{option series} is the set of contracts with the same
underlying, type, strike and \omterm{def:m1:option-contracts-and-exchanges:option}{expiry}: the unit that has a symbol, a quote and
an order book. A call is \emph{in the money}\index{in the money} when the
underlying is above its strike, a put when it is below; the amount is the
option's \emph{intrinsic value}.
\end{definition}

\begin{definition}[American and European exercise]\label{def:m1:option-contracts-and-exchanges:style}
An option with \emph{American exercise}\index{American exercise} can be
exercised on any business day up to \omterm{def:m1:option-contracts-and-exchanges:option}{expiry}; one with \emph{European
exercise}\index{European exercise} only at \omterm{def:m1:option-contracts-and-exchanges:option}{expiry}. US options on shares and
funds are American and settle by delivery of the shares; the large index
options are European and settle in cash.
\end{definition}

\begin{omfigure}
\begin{tikzpicture}
\begin{groupplot}[group style={group size=2 by 1,horizontal sep=1.6cm},
  width=6.4cm,height=4.6cm,
  xlabel={share price at expiry},
  tick label style={font=\small},label style={font=\small},
  xmin=40,xmax=60,ymin=-9,ymax=9,grid=major,grid style={gray!25},
  legend style={font=\footnotesize,at={(0.5,-0.36)},anchor=north,draw=none,fill=none,legend columns=2,/tikz/every even column/.append style={column sep=6pt}},
  table/col sep=comma]
\nextgroupplot[ylabel={P\&L per share (\$)},title={\small call, strike 50, premium 2.40}]
\addplot[omDef,very thick] table[x=s,y=long_call]{figdata/markets-1/23-option-contracts-and-exchanges/payoffs.csv};
\addplot[omThm,very thick,dashed] table[x=s,y=short_call]{figdata/markets-1/23-option-contracts-and-exchanges/payoffs.csv};
\legend{buyer,writer}
\nextgroupplot[title={\small put, strike 50, premium 1.90}]
\addplot[omDef,very thick] table[x=s,y=long_put]{figdata/markets-1/23-option-contracts-and-exchanges/payoffs.csv};
\addplot[omThm,very thick,dashed] table[x=s,y=short_put]{figdata/markets-1/23-option-contracts-and-exchanges/payoffs.csv};
\legend{buyer,writer}
\end{groupplot}
\end{tikzpicture}

\omcaption{P\&L at \omterm{def:m1:option-contracts-and-exchanges:option}{expiry}, premium included. The buyer's loss is limited to
the premium; the writer's gain is. The two lines of each panel add to zero:
options redistribute, they do not create. Data: the chapter's
build.}\label{fig:m1:option-contracts-and-exchanges:payoffs}
\end{omfigure}

A US option symbol spells its series out. Since 2010 every listed option
carries a 21-character key: six characters for the root, padded with
spaces; the \omterm{def:m1:option-contracts-and-exchanges:option}{expiry} as six digits, year, month, day; \texttt{C} or \texttt{P};
and the strike in eight digits, five for dollars and three for decimals.
\texttt{XYZ\textvisiblespace\textvisiblespace\textvisiblespace261218C00052500}
is the XYZ call expiring on 18 December 2026 with a strike of \$52.50.

How many series a name carries follows from the exchanges' listing rules:
weekly expiries for the next several weeks, monthly expiries for the next
several months, quarterly and long-dated expiries beyond, each with strikes
at intervals that widen away from the money and out in time. Five weeklies
and six monthlies with seventy strikes each, and five distant expiries with
thirty, make $2 \times (11 \times 70 + 5 \times 30) = 1\,840$ series on one
share.

\section{One clearing house}

In the United States every listed option, whichever exchange it trades on,
is issued and guaranteed by a single clearing house owned by the exchanges.
The buyer's contract is with the clearing house, not with the writer, as for
any \omterm{def:m1:clearing-and-settlement:ccp}{central counterparty} (\cref{def:m1:clearing-and-settlement:ccp}); what is
unusual is that there is only one. A call bought on one exchange can be sold
on any other, because it is the same contract with the same counterparty:
options are \emph{fungible} across exchanges, which futures are not
(\cref{ch:m1:futures-contracts-and-exchanges}).

\begin{dated}{2026-09}{Size}\label{dat:m1:option-contracts-and-exchanges:occ}
The clearing house reported 10.38 billion contracts cleared in 2022, then a
record. It \omterm{def:m1:option-contracts-and-exchanges:assignment}{assigns} exercise notices to its clearing members at random; each
member then allocates them to its customers by its own procedure.
\end{dated}

\section{The exchanges}

Fungibility made competition between exchanges possible, and the result
resembles the equity market: more than fifteen options exchanges, owned by a
handful of groups, all listing the same series; a consolidated best bid and
offer across them; a rule against trading through a better displayed price;
and a family of pricing models, from maker-taker exchanges with \omterm{def:m1:matching-algorithms-and-implied-spreads:fifo}{price-time
priority} to exchanges that charge the maker, pay for customer orders and
allocate pro rata with privileges for designated \omterm{def:m1:what-a-trading-firm-does:market-maker}{market makers}
(\cref{ch:m1:matching-algorithms-and-implied-spreads}). Index options are the
exception: the most important are licensed exclusively to one exchange
group and trade only there. The structure, and the flow that moves through
it, are the subject of \cref{ch:m1:options-market-structure}.

In Europe the picture is the futures picture. Each exchange lists options on
its home market's shares and indices and clears them in its own clearing
house; a position opened on one exchange cannot be closed on another, and
liquidity in each underlying sits in one place. Much European volume in
large size is negotiated off the book and reported to the exchange as a
block.

\section{Exercise, assignment and expiry}

\begin{definition}[Assignment]\label{def:m1:option-contracts-and-exchanges:assignment}
When a holder exercises, the clearing house selects a clearing member with a
short position in the series and \emph{assigns}\index{assignment} it the
obligation; the member passes the assignment to one of its customers who is
short. The assigned writer of a call must deliver the shares at the strike;
the assigned writer of a put must buy them. A writer learns of an
assignment the next morning.
\end{definition}

\begin{dated}{2026-09}{The rules of expiry}\label{dat:m1:option-contracts-and-exchanges:expiry}
\emph{Exercise by exception.} At \omterm{def:m1:option-contracts-and-exchanges:option}{expiry} the clearing house exercises every
option that is \omterm{def:m1:option-contracts-and-exchanges:series}{in the money} by \$0.01 or more, for equity and index options
alike, unless the clearing member instructs otherwise. It is a convenience
between the clearing house and its members, and brokers may apply their own
thresholds. \emph{Instructions.} Holders can also exercise options that are
not \omterm{def:m1:option-contracts-and-exchanges:series}{in the money}, or decline to exercise options that are: the exchanges'
cut-off for exercise notices is 16:30 Chicago time, ninety minutes after the
close, and most brokers' deadlines are earlier. \emph{Index options.}
Standard S\&P 500 index options stop trading on the Thursday and settle on
the third Friday against a value computed from the components' opening
prices; the weekly and end-of-month series trade until 16:00 New York time
on their \omterm{def:m1:option-contracts-and-exchanges:option}{expiry} date and settle against the components' closing prices.
Both are European and pay the in-the-money amount times \$100 in cash.
\end{dated}

\begin{omfigure}
\begin{tikzpicture}[font=\small,
  box/.style={draw=omDef,rounded corners=2pt,minimum width=2.5cm,minimum height=1.0cm,align=center,fill=omDef!4},
  flow/.style={-{Stealth[length=2mm]},thick}]
\node[box] (h) at (0,0) {Holder};
\node[box] (b1) at (3.3,0) {Holder's\\broker};
\node[box,draw=omThm,fill=omThm!5] (c) at (6.6,0) {Clearing\\house};
\node[box] (b2) at (9.9,0) {A short\\member};
\node[box] (w) at (13.0,0) {A writer};
\draw[flow,omDef] (h) -- (b1);
\draw[flow,omDef] (b1) -- (c);
\draw[flow,omThm] (c) -- (b2);
\draw[flow,omThm] (b2) -- (w);
\node[font=\footnotesize,text=omDef] at (3.3,1.0) {exercise notice, by the cut-off};
\node[font=\footnotesize,text=omThm] at (9.9,1.0) {assignment, overnight};
\node[font=\footnotesize] at (6.6,-1.0) {random};
\node[font=\footnotesize] at (9.9,-1.0) {member's own rule};
\end{tikzpicture}

\omcaption{Exercise and \omterm{def:m1:option-contracts-and-exchanges:assignment}{assignment}. The holder decides after the close; the
writer finds out the next morning. Between the two, the writer does not know
its own position.}\label{fig:m1:option-contracts-and-exchanges:assign}
\end{omfigure}

\begin{definition}[Pin risk]\label{def:m1:option-contracts-and-exchanges:pin}
\emph{Pin risk}\index{pin risk} is the uncertainty of a writer whose short
options expire with the underlying at or very near the strike: it cannot
know how many will be exercised, hence how many shares it will hold on
Monday, and cannot hedge what it does not know.
\end{definition}

Exercise is a decision made by holders with ninety minutes of hindsight. An
option out of the money by two cents at 16:00 is \omterm{def:m1:option-contracts-and-exchanges:series}{in the money} if the share
trades three cents higher in the after-hours session on news released at
16:05, and its holder may exercise. \Cref{fig:m1:option-contracts-and-exchanges:expiry}
shows what one \omterm{def:m1:option-contracts-and-exchanges:option}{expiry} turns a small book into, as a function of the closing
price alone; the jumps are at the strikes, and on each jump the true answer
is a random variable.

\begin{omfigure}
\begin{tikzpicture}
\begin{axis}[width=11cm,height=5cm,
  xlabel={closing price on expiry day},ylabel={net shares},
  tick label style={font=\small},label style={font=\small},
  xmin=46,xmax=54.5,ymin=-5200,ymax=7200,grid=major,grid style={gray!25},
  scaled y ticks=false,yticklabel style={/pgf/number format/fixed,/pgf/number format/1000 sep={\,}},
  table/col sep=comma]
\addplot[omDef,very thick,const plot] table[x=close,y=shares]{figdata/markets-1/23-option-contracts-and-exchanges/expiry.csv};
\node[font=\footnotesize,anchor=west] at (axis cs:46.05,5900) {short puts assigned};
\node[font=\footnotesize,anchor=west] at (axis cs:50.1,3400) {long calls exercised};
\node[font=\footnotesize,anchor=east] at (axis cs:54.45,-4350) {short calls assigned};
\end{axis}
\end{tikzpicture}

\omcaption{A book short 40 puts at 47.50 and 10 at 50, long 25 calls at 50
and short 60 calls at 52.50, all expiring today. Shares received ($+$) or
delivered ($-$) under exercise by exception, with every in-the-money short assigned in full,
against the closing price. Data: the chapter's
build.}\label{fig:m1:option-contracts-and-exchanges:expiry}
\end{omfigure}

Early exercise of American options is rational in two cases that a desk
must anticipate every day, not only at \omterm{def:m1:option-contracts-and-exchanges:option}{expiry}: a call just before an
\omterm{def:m1:shares-corporate-actions-indices:dividend}{ex-dividend date}, when the dividend exceeds the option's remaining time
value; and a deep in-the-money put, when the interest on the strike exceeds
it. The writer of such options is assigned overnight and wakes up with a
share position and without the dividend.

\section{Tutorial: symbols, a chain, payoffs}\label{tut:m1:option-contracts-and-exchanges:chain}

\begin{tutorial}
\textbf{Goal.} Parse and produce OSI symbols, hold a chain in a structure
keyed by integers, compute payoffs, and turn a book into shares at \omterm{def:m1:option-contracts-and-exchanges:option}{expiry}.
\textbf{End state:} the two data figures of this chapter.
\begin{tutsteps}
\item \textbf{A series} with its symbol, intrinsic value and payoff.
\omcode{code/firm/chain/firm_chain.py}{11}{35}{One option series. The strike is an integer of thousandths of a dollar.}\label{lst:m1:option-contracts-and-exchanges:series}
\item \textbf{Parsing.} Twenty-one characters, four fields, and an error for
anything else.
\omcode{code/firm/chain/firm_chain.py}{38}{45}{From an OSI symbol to a series.}\label{lst:m1:option-contracts-and-exchanges:parse}
\item \textbf{\omterm{def:m1:option-contracts-and-exchanges:option}{Expiry}.}
\omcode{code/firm/chain/firm_chain.py}{82}{98}{Shares received or delivered at expiry under exercise by exception.}\label{lst:m1:option-contracts-and-exchanges:expiry}
\end{tutsteps}
\textbf{What to change next.} Make \omterm{def:m1:option-contracts-and-exchanges:assignment}{assignment} random: each short contract
within five cents of the money is exercised against with a probability that
depends on the distance and on an after-hours move. Report the distribution
of Monday's share position, not a number.
\end{tutorial}

\section{Build: the option chain}\label{bld:m1:option-contracts-and-exchanges:chain}

\begin{build}
\bfield{Purpose} Every options component of the miniature firm, from the
volatility surface of One Quant Book 5 to the \omterm{def:m1:what-a-trading-firm-does:market-maker}{market maker} of Book 11, asks
one container for the series of an underlying.
\bfield{Interface} \texttt{Series(root, expiry, right, strike\_milli,
multiplier, style, settlement)} with \texttt{osi}, \texttt{strike},
\texttt{intrinsic}, \texttt{payoff}; \texttt{parse\_osi(symbol)};
\texttt{Chain(series)} with \texttt{expiries()}, \texttt{strikes(expiry)},
\texttt{get(expiry, strike\_milli, right)}, \texttt{atm\_strike(expiry,
underlying)}; \texttt{exercised\_by\_exception(series, close)};
\texttt{expiry\_shares(positions, close, overrides)}.
\bfield{Rules} Strikes are integers of thousandths everywhere; a float
strike is never a key. Duplicate series are an error. The at-the-money
strike resolves ties downwards, deterministically. \omterm{def:m1:basis-roll-and-delivery:delivery}{Cash-settled} series
deliver no shares. Corporate actions that change the deliverable
(\cref{bld:m1:shares-corporate-actions-indices:adjuster}) create adjusted
series with non-standard multipliers: the multiplier is a field, not a
constant.
\bfield{Acceptance tests} \texttt{code/firm/chain/tests/}: symbol round
trips and malformed symbols; chain lookup and the at-the-money rule; payoffs
of long and short positions; exercise by exception at one cent; net shares
of a book at several closing prices, with a contrary instruction.
\bfield{Stretch} Adjusted series after a split and a special dividend, with
the deliverable as a basket.
\end{build}

\begin{omsources}
\item The Options Industry Council, \emph{Options Exercise} (frequently asked
questions: thresholds, cut-off, assignment).
\item The Options Clearing Corporation, press release on 2022 cleared
volume, 3 January 2023; \emph{Options Symbology Initiative} documentation.
\item Cboe, \emph{S\&P 500 Index Options: product specifications} and
\emph{S\&P 500 Weeklys Options: specifications}.
\end{omsources}

\section{Exercises}

\begin{exercise}[$\star$]\label{exo:m1:option-contracts-and-exchanges:1}
Decode \texttt{XYZ\textvisiblespace\textvisiblespace\textvisiblespace270115P00047500}
and write the symbol of the XYZ call of 19 March 2027 with a strike of
\$112.50.
\end{exercise}

\begin{exercise}[$\star$]\label{exo:m1:option-contracts-and-exchanges:2}
A trader buys 10 calls with a strike of 50 for a premium of 2.40. Give the
P\&L at \omterm{def:m1:option-contracts-and-exchanges:option}{expiry} with the share at 56 and at 48, and the break-even share
price.
\end{exercise}

\begin{exercise}[$\star$]\label{exo:m1:option-contracts-and-exchanges:3}
With five weekly and six monthly expiries of seventy strikes each and five
long-dated expiries of thirty, count the series. How many quote updates does
one one-cent move of the share require from a \omterm{def:m1:what-a-trading-firm-does:market-maker}{market maker} quoting all of
them on sixteen exchanges?
\end{exercise}

\begin{exercise}[$\star\star$]\label{exo:m1:option-contracts-and-exchanges:4}
For the book of \cref{fig:m1:option-contracts-and-exchanges:expiry} give the
shares received or delivered for closing prices of 47.00, 49.00, 50.00,
50.01, 52.50 and 52.51.
\end{exercise}

\begin{exercise}[$\star\star$]\label{exo:m1:option-contracts-and-exchanges:5}
A call with a strike of 40 has 0.15 of time value left; the share trades at
46 and goes ex-dividend tomorrow with a dividend of 0.60. Should the holder
exercise today? What happens to a writer who has hedged with shares and is
assigned?
\end{exercise}

\begin{exercise}[$\star\star$]\label{exo:m1:option-contracts-and-exchanges:6}
An index \omterm{def:m1:option-contracts-and-exchanges:series}{option series} stops trading on Thursday and settles on Friday
against opening prices; another trades until Friday's close and settles
against closing prices. A desk is short both, with the same strike. Describe
its risk between Thursday's close and Friday's close in each.
\end{exercise}

\begin{exercise}[$\star\star\star$]\label{exo:m1:option-contracts-and-exchanges:7}
\emph{Coding.} With \texttt{expiry\_shares} and the chapter's book, add a
contrary instruction: the holder of the long 50 calls declines to exercise at
a close of 50.01. Give the net shares with and without it, and explain when
declining is right.
\end{exercise}

\begin{exercise}[$\star\star\star$]\label{exo:m1:option-contracts-and-exchanges:8}
\emph{Find the flaw.} ``The stock closed at 52.48. My short 52.50 calls
expired worthless, I keep the premium, and my hedge can come off at Monday's
open.'' What has the writer assumed, and what should have been done at
15:59?
\end{exercise}

\section{Problem: Expiry Friday}

\begin{problem}[{Weekend problem --- what the book is on Monday}]\label{pb:m1:option-contracts-and-exchanges:1}
A desk holds the book of \cref{fig:m1:option-contracts-and-exchanges:expiry}
in XYZ, all expiring today: short 40 puts at 47.50, short 10 puts at 50, long
25 calls at 50, short 60 calls at 52.50. Multiplier 100. It is hedged with a
short position of 1\,900 shares. XYZ closes at 52.48.

\medskip\noindent\textbf{Part I --- At the close.}
\begin{enumerate}
\item Which series are \omterm{def:m1:option-contracts-and-exchanges:series}{in the money} at 52.48, and by how much?
\item Under exercise by exception, what happens to each position?
\item Give the shares received from exercise.
\item Give the desk's share position on Monday if nothing else happens.
\item What should the desk do with that position, and when?
\end{enumerate}

\medskip\noindent\textbf{Part II --- After the close.}
At 16:10 XYZ announces a contract and trades at 53.20 in the after-hours
session.
\begin{enumerate}[resume]
\item Until when can holders of the 52.50 calls submit exercise
instructions?
\item Why would a holder exercise a call that closed out of the money?
\item On Saturday the desk learns that 45 of its 60 short calls were
assigned. Give the shares delivered.
\item Give the desk's share position on Monday.
\item XYZ opens on Monday at 54.10. Give the loss on the unexpected part of
the position, relative to the strike.
\end{enumerate}

\medskip\noindent\textbf{Part III --- What could have been done.}
\begin{enumerate}[resume]
\item At 15:55, with XYZ at 52.45, the 52.50 calls were offered at 0.04. What
would buying back the 60 have cost?
\item Compare with question 10.
\item Why does the option still cost 0.04 five minutes before it
``expires worthless''?
\item The desk could instead have bought 3\,000 shares at the close as a
precaution. Evaluate.
\item Why is the long 50 call not a worry of the same kind?
\end{enumerate}

\medskip\noindent\textbf{Part IV --- Judgement.}
\begin{enumerate}[resume]
\item \omterm{def:m1:what-a-trading-firm-does:market-maker}{Market makers}' hedging is said to ``pin'' stocks to strikes on \omterm{def:m1:option-contracts-and-exchanges:option}{expiry}
day. Give the mechanism in two sentences.
\item Would this problem exist for a \omterm{def:m1:basis-roll-and-delivery:delivery}{cash-settled} European index option?
\item A retail client is assigned on a short call in a stock that then gaps
up. Where does the loss sit until Monday?
\item State the \emph{named result}: the desk's share position on Monday
after exercise and \omterm{def:m1:option-contracts-and-exchanges:assignment}{assignment}.
\item In one sentence: who decides when an option position becomes a share
position?
\end{enumerate}
\end{problem}

\section{Interview questions}

\begin{interviewq}[$\star$ \iqroles{trader, researcher, developer}]\label{iq:m1:option-contracts-and-exchanges:1}
What is the difference between American and European options, and which
listed products are which?
\end{interviewq}

\begin{interviewq}[$\star$ \iqroles{trader, developer}]\label{iq:m1:option-contracts-and-exchanges:2}
Why can a US equity option bought on one exchange be sold on another, when a
future cannot?
\end{interviewq}

\begin{interviewq}[$\star\star$ \iqroles{trader, researcher}]\label{iq:m1:option-contracts-and-exchanges:3}
When is it rational to exercise an American call early? A put?
\end{interviewq}

\begin{interviewq}[$\star\star$ \iqroles{trader}]\label{iq:m1:option-contracts-and-exchanges:4}
What is \omterm{def:m1:option-contracts-and-exchanges:pin}{pin risk} and how do you manage it?
\end{interviewq}

\begin{interviewq}[$\star\star$ \iqroles{developer, researcher}]\label{iq:m1:option-contracts-and-exchanges:5}
Design the key for an option instrument in a trading system.
\end{interviewq}

\begin{interviewq}[$\star\star\star$ \iqroles{trader, researcher}]\label{iq:m1:option-contracts-and-exchanges:6}
You are short an index option that settles on the opening prices of \omterm{def:m1:option-contracts-and-exchanges:option}{expiry}
morning. How do you hedge the last night?
\end{interviewq}
